\section{Evaluation}
\label{sec:eval}

\subsection{Experimental Setup and Evidence Scope}
We evaluate RoboDSE in three stages: controlled execution-organization ablations,
hardware-resource exploration, and bounded projection onto an imported task
success grid. These stages distinguish hardware timing predictions from measured
closed-loop task outcomes. The frozen V7 delivery contains 39,928 retained
candidate identities and evaluates 22 latency/refresh constraint boxes.
Selected timing identities undergo 960-request checks with 192 warmup requests.
These checks establish scheduler consistency, not independent physical timing
accuracy or closed-loop validation. Conditional memory configurations, estimated
compute shapes, and static SRAM assumptions retain their documented limitations.

The imported task table contains 33 aggregate $(L,S)$ observations. We use
piecewise-linear interpolation only inside their convex hull and abstain outside
it. The table lacks per-speed trial denominators and verified checkpoint,
environment, and delay-injection provenance; its association with
$\pi_{0.5}$ and $H=16$ is an author attribution pending metadata verification.
Accordingly, projected success values are descriptive estimates, not new task
measurements, confidence bounds, or validated hardware recommendations.

\subsection{Benefits and Limits of Cross-Chunk Batching}
We hold physical hardware and the VLM/DiT core allocation fixed and compare
$k=1$, batched $k=2,3$, and unbatched $k=2,3$ organizations.
On the four-core, 102.4\,GB/s LPDDR5 configuration, unbatched execution transfers
6.407\,GB of DiT data per request. Batching reduces this to 3.293\,GB at $k=2$
and 2.670\,GB at $k=3$. Total DRAM traffic falls from 11.188\,GB/request to
8.074 and 7.451\,GB/request, respectively. With a fixed same-type DRAM energy
coefficient, these are also dynamic DRAM energy ratios, not whole-chip power
reductions.

The $k=1$ pipeline control achieves $(L,S)=(120.55,120.54)$\,ms. Batched $k=2$
and $k=3$ achieve $(178.16,89.07)$ and $(248.48,82.85)$\,ms, exchanging longer
latency for more frequent updates. This $k=1$ control is not the independently
optimized serial baseline. At matched $k$, unbatched latencies are 241.09 and
361.63\,ms while the refresh interval remains 120.54\,ms. High-bandwidth,
VLM-limited controls must also be shown: reduced traffic need not improve the
refresh interval when memory stalls are no longer limiting.

\subsection{Task-Aware Hardware-Resource Frontiers}
Across the 22 timing boxes, task-aware selection yields five strictly better
resource frontiers, three extensions, and fourteen identical frontiers. Resource
equality is defined on core-count/peak-bandwidth vector sets; ties in mapping and
storage identity are retained. For Box$(200,100)$, a four-core, 102.4\,GB/s
batched design becomes feasible alongside latency-only alternatives requiring
three cores at 204.8\,GB/s or two cores at 409.6\,GB/s. This is a bandwidth/core
tradeoff and frontier extension, not strict dominance or matched-resource success
improvement. Timing-box membership is not an empirical equal-success contour.

Area and uncertainty are secondary analyses. Under the central literature area
coefficients, zero of the 22 boxes has a strictly smaller globally selected
task-aware area. Strict-improvement task/box memberships reaching the 80\%
empirical stress-feasibility threshold number zero of 344. These findings limit
claims of robust cost superiority and motivate bandwidth margin and additional
calibration. Literature area estimates must not be described as measured silicon
area, and cross-memory-type absolute energy comparisons are not supported.

\subsection{Bounded Task-Space Projection}
Of the 39,928 retained identities, 17,992 lie inside the imported grid's convex
hull and 21,936 are excluded. For the fixed-allocation LPDDR5 controls above,
interpolation gives mean-success estimates of 46.76\%, 47.28\%, and 43.96\%
for $k=1$, batched $k=2$, and batched $k=3$, respectively. The approximately
0.52-percentage-point $k=2$ increase is a conditional interpolation diagnostic,
not a measured task gain. In particular, the 49.85\% observation at
$(200,100)$\,ms cannot be assigned to the $(178.16,89.07)$\,ms hardware point.

We choose each exact hardware identity's minimum-latency serial baseline before
consulting the task grid. In the 264 generated refreshed-memory comparisons with
covered task-aware alternatives, these baselines all lie outside the measured
hull. We therefore report no same-hardware DSE success gain. Replacing them with
slower covered serial points would introduce selection bias. Low-latency task
measurements and matched-protocol direct rollouts are required to close this gap.

% Final figures: controlled batching timing/traffic; resource Pareto examples;
% grid hull with attainable trajectories; coverage/abstention; secondary cost.
% No new real task episodes were run for this artifact.
